% !TEX program = xelatex
% 慧医数字医疗应用系统 -- 个人实习报告（李梦杰）
% 编译：xelatex 两次（目录/交叉引用）
\documentclass[UTF8,a4paper,12pt,openany,fontset=windows]{ctexrep}

\usepackage[left=3.0cm,right=2.5cm,top=2.6cm,bottom=2.6cm,headheight=15pt]{geometry}
\usepackage{fontspec}
\usepackage{xeCJK}
\usepackage{graphicx}
\usepackage{booktabs}
\usepackage{tabularx}
\usepackage{longtable}
\usepackage{array}
\usepackage{multirow}
\usepackage{float}
\usepackage{caption}
\usepackage{subcaption}
\usepackage{setspace}
\usepackage{fancyhdr}
\usepackage{enumitem}
\usepackage{amsmath}
\usepackage{amssymb}
\usepackage{hyperref}
\usepackage{xcolor}
\usepackage{microtype}

\setmainfont{Times New Roman}
\setsansfont{Arial}
\setmonofont{Consolas}
\setCJKmainfont{SimSun}
\setCJKsansfont{SimHei}
\setCJKmonofont{FangSong}

\definecolor{thesisblack}{RGB}{0,0,0}
\color{thesisblack}
\hypersetup{
  colorlinks=false,
  pdfborder={0 0 0},
  pdftitle={慧医数字医疗应用系统个人实习报告},
  pdfauthor={李梦杰}
}

\ctexset{
  chapter={
    format=\centering\heiti\zihao{-2}\bfseries\color{thesisblack},
    name={第,章},
    number=\arabic{chapter},
    beforeskip=8pt,
    afterskip=24pt,
    fixskip=true
  },
  section={
    format=\heiti\zihao{3}\bfseries\color{thesisblack},
    beforeskip=18pt,
    afterskip=10pt
  },
  subsection={
    format=\heiti\zihao{4}\bfseries\color{thesisblack},
    beforeskip=12pt,
    afterskip=6pt
  },
  contentsname={目\quad 录}
}

\setstretch{1.5}
\setlength{\emergencystretch}{3em}
\setlength{\parindent}{2em}
\setlength{\parskip}{0pt}
\setlist{nosep,leftmargin=2em}
\captionsetup{font=small,labelfont=bf,textfont=normalfont,labelsep=quad}
\renewcommand{\arraystretch}{1.35}
% 图片路径：优先本目录图片，其次组报告图片目录与证据目录
\graphicspath{{图片/}{../最终交付文档/图片/}{../process-docs/evidence/}}

\pagestyle{fancy}
\fancyhf{}
\fancyhead[C]{\songti\zihao{5}慧医数字医疗应用系统个人实习报告}
\fancyfoot[C]{\zihao{5}\thepage}
\renewcommand{\headrulewidth}{0.4pt}
\renewcommand{\footrulewidth}{0pt}
\fancypagestyle{plain}{
  \fancyhf{}
  \fancyhead[C]{\songti\zihao{5}慧医数字医疗应用系统个人实习报告}
  \fancyfoot[C]{\zihao{5}\thepage}
  \renewcommand{\headrulewidth}{0.4pt}
}

\newcolumntype{Y}{>{\centering\arraybackslash}X}
\newcolumntype{L}{>{\raggedright\arraybackslash}X}
\newcommand{\thesisfigure}[3][0.96\textwidth]{%
  \begin{figure}[H]
    \centering
    \includegraphics[width=#1,height=0.72\textheight,keepaspectratio]{#2}
    \caption{#3}
  \end{figure}
}
\newcommand{\codefont}{\ttfamily\footnotesize}

\begin{document}
\songti\zihao{-4}
\hypersetup{pageanchor=false}

% ============ 封面 ============
\begin{titlepage}
  \thispagestyle{empty}
  \begin{center}
    \vspace*{0.8cm}
    {\heiti\zihao{1}\bfseries 个人实习报告\par}
    \vspace{1.6cm}
    {\heiti\zihao{2}\bfseries 慧医数字医疗应用系统\par}
    \vspace{0.6cm}
    {\songti\zihao{3}软件工程综合实训项目\par}
    \vfill
    \begin{minipage}{0.78\textwidth}
      \zihao{4}
      \noindent\heiti 学\quad 校：\songti 西南民族大学\par\vspace{0.55em}
      \noindent\heiti 姓\quad 名：\songti 李梦杰\par\vspace{0.55em}
      \noindent\heiti 专\quad 业：\songti 软件工程\par\vspace{0.55em}
      \noindent\heiti 班\quad 级：\songti 软件工程2302\par\vspace{0.55em}
      \noindent\heiti 学\quad 号：\songti 202331209061\par\vspace{0.55em}
      \noindent\heiti 项目角色：\songti 项目经理 \& 开发工程师\par\vspace{0.55em}
      \noindent\heiti 指导教师：\songti 万华军\par\vspace{0.55em}
      \noindent\heiti 完成日期：\songti 2026年7月17日\par
    \end{minipage}
    \vfill
    {\songti\zihao{4}2026年7月\par}
  \end{center}
\end{titlepage}

\hypersetup{pageanchor=true}
\pagenumbering{Roman}

% ============ 摘要 ============
\chapter*{摘\quad 要}
\addcontentsline{toc}{chapter}{摘要}
本次实训以“慧医数字医疗应用系统”为载体，在 2026 年 7 月 13 日至 7 月 16 日共四个工作日内，完成了一个前后端分离、面向医药代表业务管理的 Web 系统的需求分析、架构设计、编码实现、自动化测试与容器化部署。本人李梦杰（学号 202331209061，软件工程 2302 班）在项目中同时担任\textbf{项目经理}与\textbf{开发工程师}双重角色：作为项目经理，负责 5 人团队的角色分工与 owner 映射、Epic--Feature--Story--Task 分层工作项分解、五日排期与里程碑设定、双仓库（GitHub + 华为云 CodeArts）协作规范制定、CI/CD 流水线落地、代码检查与覆盖率门禁固化，以及过程文档与证据材料的组织归档；作为开发工程师，主导 EP-01“统一身份与权限中心”的后端实现，承担认证、权限树、httpOnly Cookie 会话、仪表盘聚合与高德逆地理编码代理等模块，并负责报告撰写与代码扫描整改。

系统后端基于 Spring Boot 2.7 与 JDK 17 实现规范化 RBAC 权限控制与会话管理，前端基于 Vue 3 构建医药代表工作台，通过 Nginx 同源代理与 Docker Compose 容器化部署，经华为云 CodeArts 流水线（\#12 全流程通过，耗时 5 分 1 秒）发布至 ECS，并上线正式子域名 \texttt{sx.weavecan.com}。本人合计完成 295 次提交（含项目初始化与 15 次功能分支集成合并），覆盖需求、设计、实现、测试、部署全流程；后端 Java 代码 6\,387 行、前端 Vue/JS 代码 9\,294 行，代码级测试 241/241 通过，后端与前端行覆盖率均达 100\%。

\textbf{关键词}：数字医疗；前后端分离；RBAC；Docker 部署；CodeArts 持续集成；项目管理；UML 建模

% ============ 目录 ============
\tableofcontents
\cleardoublepage
\pagenumbering{arabic}

% =============================================================
\chapter{实习概述}

\section{实习背景与目的}
医疗机构及医药相关组织在日常管理中需要维护医生、药品、政策、材料、企业和销售地点等多类基础数据。传统分散式表格管理容易产生数据重复、权限边界模糊、查询效率低和变更难以追踪等问题。本次软件工程综合实训以“慧医数字医疗应用系统”为载体，要求在一个紧凑周期内完成从需求分析、UML 建模、数据库设计、前后端实现、自动化测试到容器部署与持续集成的完整软件工程闭环，旨在训练真实工程环境下的系统设计能力、团队协作能力与工程交付能力。

\section{项目简介}
本项目建成一个面向医药代表业务管理的 Web 系统，以管理员与医生两类角色建立清晰的最小权限边界，覆盖医生、药品、医药公司、城市、销售地点（含地图）、医保政策、公司政策、必备材料与图片上传等八类核心业务管理，并提供首页统计面板与健康检查能力。系统以 MySQL 数据库 \texttt{medicine} 为权威数据源，以 Redis 保存会话与缓存，前端通过 Nginx 同源代理访问后端 \texttt{/api} 接口，生产环境经华为云 ECS 与正式子域名对外提供服务。

\section{个人角色与职责}
本人在项目中承担\textbf{项目经理}与\textbf{开发工程师}双重角色，对应过程文档中的账号 \texttt{anmory666}（GitHub 用户名 \texttt{anmoryli}）。具体职责见表\ref{tab:my-role}。

\begin{table}[H]
  \centering
  \caption{个人角色与职责}\label{tab:my-role}
  \begin{tabularx}{\textwidth}{p{2.6cm}L}
    \toprule
    角色 & 主要职责 \\
    \midrule
    项目经理 & 5 人团队角色分工与 owner 映射；Epic/Feature/Story/Task 分层工作项分解；五日排期与 5 个里程碑设定；双仓库协作规范制定；CodeArts 流水线推动；代码检查与覆盖率门禁固化；过程文档与证据归档 \\
    开发工程师 & EP-01 统一身份与权限中心（登录与会话、权限树、httpOnly Cookie）；仪表盘聚合接口；高德逆地理编码代理；图片上传安全；主导 httpOnly Cookie 持久登录改造 \\
    系统开发工程师 & 系统整体架构与技术选型；前后端契约与联调；数据库迁移；部署拓扑；UML 设计文档；报告撰写与代码扫描整改 \\
    集成者 & 审核并合并 15 个功能分支 PR 至主分支，保证双仓库同步与历史一致 \\
    \bottomrule
  \end{tabularx}
\end{table}

\section{实习时间与成果概览}
实习开发周期为 2026 年 7 月 13 日至 7 月 16 日，共四个工作日（计划延伸至 7 月 17 日验收）。期间本人合计完成 295 次提交。提交按日期分布见表\ref{tab:daily-commits}。

\begin{table}[H]
  \centering
  \caption{提交按日期分布}\label{tab:daily-commits}
  \begin{tabularx}{\textwidth}{p{3.0cm}p{2.5cm}L}
    \toprule
    日期 & 提交数 & 主要工作 \\
    \midrule
    2026-07-13 & 25 & 项目初始化、需求基线、数据库 V1--V2、后端脚手架、57 项 API 回归 57/57 \\
    2026-07-14 & 88 & CI/CD 流水线、双仓库规范、Vue3+Vite 迁移、httpOnly Cookie、UML 渲染 \\
    2026-07-15 & 180 & 测试补全与 100\% 行覆盖、代码扫描整改、文档归档 \\
    2026-07-16 & 2 & 报告表格修复与代码检查 G.FMT.03 整改 \\
    \midrule
    合计 & 295 & 覆盖需求、设计、实现、测试、部署全流程 \\
    \bottomrule
  \end{tabularx}
\end{table}

\section{项目小组信息}
\begin{table}[H]
  \centering
  \caption{项目小组信息}
  \begin{tabularx}{\textwidth}{p{3.2cm}L}
    \toprule
    项目项 & 内容 \\
    \midrule
    学校/专业/班级 & 西南民族大学 / 软件工程 / 软件工程2302 \\
    小组名称 & 02组 \\
    小组组长 & 李梦杰（项目经理 \& 开发） \\
    小组成员 & 曾子唯、詹宁康、周林桂、侬常旭 \\
    指导教师 & 万华军 \\
    完成时间 & 2026年7月17日 \\
    \bottomrule
  \end{tabularx}
\end{table}

% =============================================================
\chapter{项目管理与进度安排}

作为项目经理，本人在项目启动阶段即完成了团队角色分工、工作项分解、排期里程碑与协作规范的整体设计，并形成可追溯的过程文档沉淀在 \texttt{process-docs/} 目录。本章从项目经理视角说明这些安排如何落地。

\section{团队组建与角色分工}
项目组共 5 人，按“项目管理 + 开发 + CI/CD + 测试”划分角色，每个角色在 \texttt{process-docs/} 下有独立的职责说明文档，且每个 Epic/Feature/Story/Task 的负责人字段直接写入真实账号，统一录入华为云 ProjectMan。角色分工见表\ref{tab:roles}。

\begin{table}[H]
  \centering
  \caption{项目角色分工}\label{tab:roles}
  \begin{tabularx}{\textwidth}{p{3.0cm}p{2.8cm}L}
    \toprule
    角色 & 账号 & 职责 \\
    \midrule
    项目经理/开发 & anmory666（本人） & 工作项分解、里程碑、规范制定；认证与权限中心后端开发 \\
    CI/CD 工程师 & Ning\_24 & 流水线搭建、双仓库同步、部署健康检查、基础资料（城市/公司/销售地点） \\
    开发工程师 & khyxone & 医生、药品业务模块 \\
    开发工程师 & spikezzw & 必备材料、医药公司政策、医保政策、图片上传 \\
    测试工程师 & mengyetianxing & 前端联调、测试交付、黑盒回归、覆盖率门禁 \\
    \bottomrule
  \end{tabularx}
\end{table}

\section{工作项分解}
本人采用 \textbf{Epic--Feature--Story--Task} 四层结构组织工作项，并以 CSV 落地、录入 ProjectMan，形成可追溯的层级关系（Task 通过“父工作项”字段挂到对应 Story）。分解规模见表\ref{tab:backlog}。

\begin{table}[H]
  \centering
  \caption{工作项分解规模}\label{tab:backlog}
  \begin{tabularx}{\textwidth}{p{2.4cm}p{1.6cm}L}
    \toprule
    层级 & 数量 & 说明 \\
    \midrule
    Epic & 5 & EP-01 身份与权限中心、EP-02 业务基础资料、EP-03 医生与药品、EP-04 政策材料与文件服务、EP-05 交付与质量保障 \\
    Feature & 7 & 在 Epic 下进一步拆分的功能特性 \\
    Story & 15 & 含工时估算（6h--12h），如登录与会话管理、权限树与角色边界、仪表盘聚合等 \\
    Task & 57 & 黑盒验收任务，每个 1h，对应 57 个接口场景 \\
    里程碑 & 5 & MS-01--MS-05，每个含验收标准与负责人 \\
    \bottomrule
  \end{tabularx}
\end{table}

\section{进度排期与里程碑}
计划周期为 2026-07-13 至 2026-07-17，共 5 个工作日。本人设定 5 个里程碑，每个里程碑都有明确验收标准与负责人，见表\ref{tab:milestones}。其中 MS-01（需求与任务基线）与 MS-02（核心后端功能完成）由本人负责。

\begin{table}[H]
  \centering
  \caption{里程碑计划}\label{tab:milestones}
  \begin{tabularx}{\textwidth}{p{1.4cm}p{3.2cm}p{2.8cm}L}
    \toprule
    里程碑 & 名称 & 负责人/日期 & 验收标准 \\
    \midrule
    MS-01 & 需求与任务基线 & anmory666 / 07-13 & Epic、Feature、Story、Task 层级确认 \\
    MS-02 & 核心后端功能完成 & anmory666 / 07-15 & 认证、仪表盘、基础资料接口可用 \\
    MS-03 & 业务模块完成 & khyxone / 07-16 & 医生、药品、材料、政策和上传可用 \\
    MS-04 & 前端联调完成 & mengyetianxing / 07-17 & 主要页面可访问且接口联调通过 \\
    MS-05 & 发布验收完成 & mengyetianxing / 07-17 & 57 条黑盒场景通过、Docker 健康、文档齐全 \\
    \bottomrule
  \end{tabularx}
\end{table}

Epic 级排期按负责人与日期分配：EP-01 由本人在 07-13--07-14 完成（P0 关键）；EP-02 由 Ning\_24 在 07-13--07-15；EP-03 由 khyxone 在 07-14--07-16；EP-04 由 spikezzw 在 07-13--07-17；EP-05 由 mengyetianxing 贯穿全程。

\section{双仓库协作规范}
本人撰写并落地 \texttt{双仓库提交规范.md}，作为团队长期有效的协作基线。项目同时使用两个远程仓库：\texttt{origin} 指向 GitHub（\texttt{anmoryli/shixun-swmu}，默认分支 \texttt{main}），\texttt{codearts} 指向华为云 CodeArts（SSH，默认分支 \texttt{master}）。本地分支统一为 \texttt{main}，通过 \texttt{main:master} refspec 推到 CodeArts，每次提交依次执行：
\begin{table}[H]
  \centering
  \caption{核心代码：双仓库推送命令}
  \begin{tabularx}{\textwidth}{p{1.2cm}L}
    \toprule
    行 & 代码 \\
    \midrule
    1 & \codefont git push origin main            \# GitHub: main -> origin/main \\
    2 & \codefont git push codearts main:master   \# CodeArts: main -> codearts/master \\
    \bottomrule
  \end{tabularx}
\end{table}

规范同时要求：提交粒度最小化（写完一个最小功能单元即提交，如一个函数、一条扫描规则、一段文档）；提交信息格式 \texttt{<类型>(<范围>): <描述>}，类型含 feat/fix/refactor/style/docs/test/chore；推送前校验后端 \texttt{mvn -o clean compile}、前端 \texttt{npm run build}、Docker 改动执行 \texttt{docker compose up -d --build}；日常提交均为 fast-forward，禁止强推覆盖他人历史。

\thesisfigure{GitHub仓库.png}{GitHub 项目仓库}
\thesisfigure{GitHubPR.png}{GitHub Pull Request 协作记录}
\thesisfigure{华为云代码仓库1.png}{CodeArts 代码仓库概览（一）}

\section{CodeArts 流水线与持续集成}
本人配合 CI/CD 工程师配置 CodeArts 流水线 \texttt{anmory-20260714110641}，源为 CodeArts Repo \texttt{shixun-2/master}，区域华北-北京四。最终流水线 \#12 对提交 \texttt{41ade53e} 全流程通过，总耗时 5 分 1 秒，部署目标为 ECS \texttt{1.14.150.130}，演示端口 9092。阶段结构见表\ref{tab:pipeline}。

\begin{table}[H]
  \centering
  \caption{CodeArts 流水线阶段}\label{tab:pipeline}
  \begin{tabularx}{\textwidth}{p{2.8cm}p{3.6cm}p{2.2cm}L}
    \toprule
    阶段 & 任务 & 结果 & 耗时 \\
    \midrule
    构建和检查 & 构建（JDK 17 + Maven + Node） & 成功 & 2 分 22 秒 \\
    构建和检查 & 代码检查（Java/JS/Python） & 成功 & 1 分 21 秒 \\
    部署和测试 & 部署到 shixun-prod & 成功 & 1 分 41 秒 \\
    接口测试 & Shell 健康检查 /actuator/health & 成功 & 56 秒 \\
    \bottomrule
  \end{tabularx}
\end{table}

日常协作流程为：成员在功能分支开发并提 PR $\rightarrow$ 评审通过后合并 \texttt{master} $\rightarrow$ CodeArts 监听 \texttt{master} 提交并启动流水线 $\rightarrow$ 构建与代码检查必须全成功 $\rightarrow$ 制品部署到 \texttt{shixun-prod} $\rightarrow$ 公网健康检查成功则发布完成。本人作为集成者，审核并合并了全部 15 个功能分支 PR。

\thesisfigure{华为CICD.png}{CodeArts 持续集成与持续交付概览}
\thesisfigure{华为云编译构建.png}{CodeArts 编译构建任务}
\thesisfigure{华为云部署.png}{CodeArts 部署任务}
\thesisfigure{华为云仪表盘.png}{CodeArts 项目仪表盘}

\section{代码检查与质量门禁}
项目通过 4 轮华为云 CodeArts 代码扫描报告（07-13 至 07-15）驱动 finding 修复，前后端 G.* 系列规则全部闭环。本人同时推动质量门禁固化：后端 \texttt{pom.xml} 接入 JaCoCo 0.8.12，\texttt{mvn verify} 在 bundle 行覆盖率低于 1.00 时直接失败；前端 \texttt{vite.config.js} 设置 statements/lines 100\% 阈值，低于 100\% 退出失败。这使得覆盖率从初期的 46.89\% 行覆盖提升到 100\% 并被构建硬门禁锁定。

\thesisfigure{华为云代码检查.png}{CodeArts 代码检查结果}

\section{过程证据归档机制}
本人建立 \texttt{process-docs/evidence/} 三级目录归档证据，同时保存图片与结构化文本，与各角色 README 交叉引用，确保交付可验证、可复核。证据分类见表\ref{tab:evidence}。

\begin{table}[H]
  \centering
  \caption{过程证据归档}\label{tab:evidence}
  \begin{tabularx}{\textwidth}{p{4.2cm}p{3.2cm}L}
    \toprule
    目录 & 内容 & 形式 \\
    \midrule
    evidence/api & 5 轮 API 回归报告（最终 57/57） & png/svg + json/md/xml \\
    evidence/deployment & 部署健康检查与 Docker Compose 拓扑 & png/svg + md \\
    evidence/requirements & 需求验收、覆盖、数据与 API、范围 & png \\
    \bottomrule
  \end{tabularx}
\end{table}

证据收集遵循可复核（图片 + 结构化文本双存）、时间戳化（文件名带 UTC 时间戳保留历史轮次）、不落盘敏感信息（凭据仅从进程环境读取，报告不含用户名/密码/Token）三项原则。

\section{个人承担的测试任务}
除开发与管理外，本人在 \texttt{task-backlog.csv} 中承担 10 个黑盒验收任务，集中在认证与权限边界：TK-03 伪造 Token 访问、TK-04 登录缺密码、TK-05 错误密码登录、TK-06 管理员 JSON 登录兼容、TK-07 JSON 登录会话退出、TK-08 管理员表单登录、TK-09 管理员权限树、TK-33 默认重置密码登录、TK-56 管理员退出登录、TK-57 退出后 Token 失效。

% =============================================================
\chapter{需求分析}

\section{用户角色分析}
系统包含管理员与医生两类主要角色。管理员具有业务基础数据的维护权限，负责新增、修改、删除和查询；医生以业务查询和使用为主，不应取得系统级维护权限。所有受保护请求必须经过会话校验和角色授权，未登录用户只能访问登录、会话探测和必要的健康检查入口。

\begin{table}[H]
  \centering
  \caption{角色与权限需求}
  \begin{tabularx}{\textwidth}{p{3.6cm}p{3.0cm}L}
    \toprule
    角色 & 权限范围 & 主要职责 \\
    \midrule
    管理员（ROLE\_ADMIN） & 全量业务维护 & 管理账号、医生、药品、企业、政策、城市、销售地点、材料和图片 \\
    医生（ROLE\_DOCTOR） & 授权范围内查询 & 查看首页统计和业务信息，使用地图与资讯，不执行越权维护 \\
    未登录用户 & 公开入口 & 登录、会话状态探测和网关允许的健康检查 \\
    \bottomrule
  \end{tabularx}
\end{table}

\section{功能需求}
系统主要功能需求见表\ref{tab:fr}，覆盖登录认证、权限控制、首页统计与八类业务管理。

\begin{longtable}{p{1.4cm}p{3.1cm}p{4.6cm}p{4.6cm}}
  \caption{主要功能需求}\label{tab:fr}\\
  \toprule
  编号 & 模块 & 功能说明 & 验收要点 \\
  \midrule
  \endfirsthead
  \toprule
  编号 & 模块 & 功能说明 & 验收要点 \\
  \midrule
  \endhead
  F01 & 登录认证 & 校验账号、密码与账号状态，建立服务端会话 & 正确凭据登录；错误凭据拒绝；Cookie 为 httpOnly \\
  F02 & 权限控制 & 根据真实角色返回权限树并控制接口访问 & 管理员与医生权限隔离；越权请求返回拒绝 \\
  F03 & 首页统计 & 聚合医生、药品、城市等数量与资讯 & 数据可视化正确；业务变更后缓存失效 \\
  F04 & 医生管理 & 医生信息分页、查询、新增、修改、删除与密码重置 & 唯一性、状态和权限校验正确 \\
  F05 & 医药公司 & 企业基础信息与公司政策维护 & 分页和增删改查完整 \\
  F06 & 销售地点 & 销售地点与经纬度维护，地图展示与选点 & 坐标可保存、回显并在地图中定位 \\
  F07 & 城市管理 & 城市信息分页和维护 & 数据校验与关联使用正常 \\
  F08 & 药品管理 & 药品信息、价格、公司关联与图片维护 & 图片路径安全，列表与详情一致 \\
  F09 & 政策管理 & 医保政策与公司政策维护 & 富文本或长文本完整保存和查询 \\
  F10 & 材料管理 & 必备材料信息维护 & 数据增删改查和权限控制正确 \\
  F11 & 健康检查 & 汇总应用、MySQL 和 Redis 状态 & 容器与流水线可以据此判定部署成功 \\
  \bottomrule
\end{longtable}

\section{非功能需求}
\begin{table}[H]
  \centering
  \caption{非功能需求}
  \begin{tabularx}{\textwidth}{p{2.6cm}L p{4.0cm}}
    \toprule
    类别 & 需求描述 & 评价方式 \\
    \midrule
    安全性 & 密码摘要存储；会话不可被脚本读取；文件名和路径受控；生产后端端口不直接暴露公网 & 安全测试、配置检查、越权用例 \\
    可靠性 & 应用异常统一处理；数据库和 Redis 状态可观测；部署失败可回滚 & 自动化测试、健康检查、流水线记录 \\
    可维护性 & 前后端分层、模块边界清楚、配置外置、数据库迁移可追踪 & 代码结构和文档审查 \\
    性能 & 列表分页；首页聚合结果缓存；静态资源由 Nginx 提供 & 接口耗时与历史 API 回归数据 \\
    可部署性 & 支持 JDK 17、Node 18、Docker Compose 与 CodeArts 流水线 & 构建、容器启动和云端部署验证 \\
    可测试性 & 核心模块具有单元或组件测试，覆盖率可量化 & Maven、Vitest/V8、Python 测试报告 \\
    \bottomrule
  \end{tabularx}
\end{table}

\section{业务边界与约束}
系统以 MySQL 数据库 \texttt{medicine} 为权威业务数据源，以 Redis 保存会话与缓存。前端不得直接连接数据库或 Redis；所有写操作通过后端接口完成。生产环境通过 Nginx 暴露 80/443 端口，后端 8082 仅在主机内部或容器网络中可达。药品图片只能写入指定上传目录，并限制扩展名、MIME 类型、大小和规范化路径。

\thesisfigure{01-用例图.png}{系统用例图}

% =============================================================
\chapter{系统设计}

\section{设计思想与总体架构}
系统采用“分层解耦、权限前置、配置外置、测试护航、容器交付”的总体思想。浏览器端负责交互与展示，Nginx 负责统一入口和反向代理，Spring Boot 负责业务规则、安全和数据访问，MySQL 负责持久化，Redis 负责会话和热点缓存。系统分层架构见表\ref{tab:arch}。

\begin{table}[H]
  \centering
  \caption{系统分层架构}\label{tab:arch}
  \begin{tabularx}{\textwidth}{p{2.6cm}p{4.0cm}L}
    \toprule
    层次 & 主要技术 & 职责 \\
    \midrule
    表示层 & Vue 3、Element Plus、ECharts、地图组件 & 页面交互、表单校验、数据展示和路由导航 \\
    接入层 & Nginx & 静态资源、同源代理、TLS 终止、请求头转发 \\
    应用层 & Spring Boot、Spring MVC、Spring Security & 接口编排、认证授权、统一响应与异常处理 \\
    领域与服务层 & Service、事件、缓存维护任务 & 业务规则、事务、数据变更事件与缓存失效 \\
    数据访问层 & MyBatis & SQL 映射、分页查询和实体持久化 \\
    基础设施层 & MySQL 8、Redis 7、Docker Compose & 持久数据、会话缓存、容器运行和健康检查 \\
    \bottomrule
  \end{tabularx}
\end{table}

\section{技术选型}
\begin{table}[H]
  \centering
  \caption{主要技术选型}
  \begin{tabularx}{\textwidth}{p{2.8cm}p{4.2cm}L}
    \toprule
    类别 & 技术与版本 & 选型说明 \\
    \midrule
    前端框架 & Vue 3.5、Vue Router 4、Vuex 4 & 组件化开发、动态路由和集中状态管理 \\
    构建与测试 & Vite 6、Vitest、V8 Coverage & 快速构建、组件测试和覆盖率统计 \\
    后端框架 & Spring Boot 2.7.18、JDK 17 & 成熟的 Web、配置、测试和健康检查生态 \\
    安全与持久化 & Spring Security、MyBatis & 统一安全链与可控 SQL 映射 \\
    数据与缓存 & MySQL 8、Redis 7 & 关系数据持久化、会话和热点数据缓存 \\
    交付平台 & Docker Compose、Nginx、CodeArts & 环境一致性、反向代理和持续交付 \\
    \bottomrule
  \end{tabularx}
\end{table}

\thesisfigure{02-领域类图.png}{领域类图}
\thesisfigure{06-组件图.png}{系统组件图}
\thesisfigure{07-包图.png}{系统包图}
\thesisfigure{08-部署图.png}{系统部署图}

\section{登录认证与规范化 RBAC 设计}
登录时后端校验账号状态和 BCrypt 密码摘要，成功后生成随机不透明令牌，仅将令牌的 SHA-256 摘要键与最小会话身份写入 Redis，并通过 httpOnly Cookie 交付浏览器。后续请求由安全过滤链提取 Cookie、校验 Redis 会话，再以会话中的可信账号编号为入口重新加载当前授权。原始令牌不写入 Redis，也不在登录响应体中暴露。

V5 迁移将原来的“账号字符串角色 + 菜单映射”升级为规范化的账号--角色--权限模型，核心表映射见表\ref{tab:rbac}。运行时授权链路为“\texttt{accountId} $\rightarrow$ \texttt{account\_role} $\rightarrow$ 启用的 \texttt{rbac\_role} $\rightarrow$ \texttt{rbac\_role\_permission} $\rightarrow$ 启用的 \texttt{permission}”。\texttt{TokenAuthenticationFilter} 从会话主体取得 \texttt{accountId}，逐次请求联查账号状态、角色和权限，将 \texttt{ROLE\_ADMIN}、\texttt{ROLE\_DOCTOR} 角色权限及 \texttt{resource:verb} 动作权限写入 Spring Security 的 \texttt{GrantedAuthority}。

\begin{table}[H]
  \centering
  \caption{规范化 RBAC 核心表映射}\label{tab:rbac}
  \begin{tabularx}{\textwidth}{p{3.2cm}p{5.1cm}L}
    \toprule
    表 & 关键字段 & 职责 \\
    \midrule
    \texttt{account} & \texttt{id}、\texttt{status}、\texttt{utype} & 认证主体与启停状态 \\
    \texttt{rbac\_role} & \texttt{id}、\texttt{code}、\texttt{enabled} & 角色主数据，启用 ADMIN 与 DOCTOR \\
    \texttt{account\_role} & \texttt{account\_id}、\texttt{role\_id}、\texttt{is\_primary} & 账号与角色多对多关系 \\
    \texttt{permission} & \texttt{code}、\texttt{permission\_type}、\texttt{enabled} & 统一权限目录；MENU/ACTION 区分 \\
    \texttt{rbac\_role\_\allowbreak permission} & \texttt{role\_id}、\texttt{permission\_id} & 角色与权限多对多授权关系 \\
    \bottomrule
  \end{tabularx}
\end{table}

前后端按“展示控制”和“安全控制”分层：\texttt{/api/permissions} 只根据 \texttt{accountId} 查询多角色权限并集返回菜单树与权限码，客户端传入 \texttt{roleName} 不参与授权；Vuex 与 \texttt{\$can()} 只负责菜单与按钮展示；后端控制器通过 \texttt{@PreAuthorize(hasAuthority(...))} 校验 \texttt{ACTION} 权限，才是不可绕过的最终安全边界。角色撤销、权限调整和账号禁用在下一次请求即可生效。

\thesisfigure{09-登录活动图.png}{登录活动图}
\thesisfigure{12-账号状态图.png}{账号状态图}
\thesisfigure{14-登录授权顺序图.png}{登录授权顺序图}

\section{数据库设计}
数据库使用 \texttt{medicine} 模式，基线脚本定义初始结构，迁移脚本逐步补充业务字段、索引、坐标、审计记录和规范化 RBAC 关系。执行 V2--V5 后共包含 23 张业务及关联表，其中 V5 在原 20 张表基础上新增 3 张 RBAC 关系表。脚本演进见表\ref{tab:sql}。

\begin{table}[H]
  \centering
  \caption{数据库脚本演进}\label{tab:sql}
  \begin{tabularx}{\textwidth}{p{3.2cm}p{3.0cm}L}
    \toprule
    脚本 & 阶段 & 主要作用 \\
    \midrule
    \texttt{medical.sql} & 基线 & 创建初始业务表、约束和基础数据 \\
    V2 迁移 & 需求对齐 & 增加账号状态、销售地点坐标、资讯、密码重置审计等 \\
    V3 迁移 & 审计生命周期 & 密码重置审计外键调整为随账号删除级联清理 \\
    V4 迁移 & 地图菜单 & 增加销售地点地图菜单并补充旧角色菜单映射 \\
    V5 迁移 & 规范化 RBAC & 新增账号--角色--权限关系、MENU/ACTION 权限码及回填 \\
    \bottomrule
  \end{tabularx}
\end{table}

\thesisfigure{04-实体关系图ER.png}{数据库实体关系图}

\section{关键设计要点}
图片上传由后端统一控制：只允许配置的图片类型和扩展名，限制文件大小，对文件名随机化，并将目标路径规范化后验证其仍位于上传根目录内。首页面板将统计数量、分布和资讯结果缓存在 Redis 中，业务服务完成数据变更后发布事件，监听器统一删除相关缓存；定时维护任务作为补充机制，缓存不可用时服务可从数据库读取降级运行。

\thesisfigure{11-药品图片上传活动图.png}{药品图片上传活动图}
\thesisfigure{18-首页面板顺序图.png}{首页面板顺序图}

% =============================================================
\chapter{系统实现}

本章选取能够说明关键设计的短代码片段，所有代码均在三线表中展示，完整实现以仓库源代码为准。本人主导的 EP-01 认证与权限中心相关代码是重点。

\section{安全认证边界}
\begin{table}[H]
  \centering
  \caption{核心代码：公共入口与认证边界}
  \begin{tabularx}{\textwidth}{p{1.2cm}L}
    \toprule
    行 & 代码 \\
    \midrule
    1 & \codefont .antMatchers(HttpMethod.OPTIONS, "/**").permitAll() \\
    2 & \codefont .antMatchers(HttpMethod.GET, "/image/**").permitAll() \\
    3 & \codefont .antMatchers("/api/login", "/api/session", \\
    4 & \codefont \quad "/actuator/health", "/actuator/info").permitAll() \\
    \bottomrule
  \end{tabularx}
\end{table}

\section{安全会话 Cookie}
\begin{table}[H]
  \centering
  \caption{核心代码：安全会话 Cookie}
  \begin{tabularx}{\textwidth}{p{1.2cm}L}
    \toprule
    行 & 代码 \\
    \midrule
    1 & \codefont ResponseCookie.ResponseCookieBuilder builder = \\
    2 & \codefont \quad ResponseCookie.from(cookieProperties.getName(), token) \\
    3 & \codefont \quad .httpOnly(true).secure(cookieProperties.isSecure()) \\
    4 & \codefont \quad .sameSite(cookieProperties.getSameSite()) \\
    5 & \codefont \quad .path(cookieProperties.getPath()); \\
    \bottomrule
  \end{tabularx}
\end{table}

\section{业务变更触发缓存失效}
\begin{table}[H]
  \centering
  \caption{核心代码：业务变更触发缓存失效}
  \begin{tabularx}{\textwidth}{p{1.2cm}L}
    \toprule
    行 & 代码 \\
    \midrule
    1 & \codefont @AfterReturning("businessMutation()") \\
    2 & \codefont public void publishDataChangedEvent(JoinPoint joinPoint) \{ \\
    3 & \codefont \quad eventPublisher.publishEvent( \\
    4 & \codefont \quad\quad new BusinessDataChangedEvent( \\
    5 & \codefont \quad\quad joinPoint.getSignature().toShortString())); \\
    6 & \codefont \} \\
    \bottomrule
  \end{tabularx}
\end{table}

\section{前端统一 API 请求}
\begin{table}[H]
  \centering
  \caption{核心代码：前端统一 API 请求}
  \begin{tabularx}{\textwidth}{p{1.2cm}L}
    \toprule
    行 & 代码 \\
    \midrule
    1 & \codefont export function getDoctorInfo(pn, size, keyword = '') \{ \\
    2 & \codefont \quad return request(\{ url: '/doctors', method: 'GET', \\
    3 & \codefont \quad\quad params: \{ pn, size, keyword \} \}) \\
    4 & \codefont \quad\quad .then((res) $=>$ judgeQueryResult(res)); \\
    5 & \codefont \} \\
    \bottomrule
  \end{tabularx}
\end{table}

\section{Nginx 同源代理}
\begin{table}[H]
  \centering
  \caption{核心代码：Nginx 同源代理}
  \begin{tabularx}{\textwidth}{p{1.2cm}L}
    \toprule
    行 & 代码 \\
    \midrule
    1 & \codefont location /api/ \{ \\
    2 & \codefont \quad proxy\_pass \$medicine\_backend\_upstream\$request\_uri; \\
    3 & \codefont \quad proxy\_http\_version 1.1; \\
    4 & \codefont \quad proxy\_set\_header Origin ""; \\
    5 & \codefont \} \\
    \bottomrule
  \end{tabularx}
\end{table}

\section{Docker Compose 健康检查}
\begin{table}[H]
  \centering
  \caption{核心代码：Docker Compose 健康检查}
  \begin{tabularx}{\textwidth}{p{1.2cm}L}
    \toprule
    行 & 代码 \\
    \midrule
    1 & \codefont healthcheck: \\
    2 & \codefont \quad test: ["CMD-SHELL", "wget -q -O - \\
    3 & \codefont \quad\quad http://127.0.0.1:8082/actuator/health \\
    4 & \codefont \quad\quad | grep -q '"status":"UP"'"] \\
    5 & \codefont \quad interval: 10s \\
    \bottomrule
  \end{tabularx}
\end{table}

\section{系统运行效果}
\thesisfigure{项目效果登录页面.png}{系统登录页面}
\thesisfigure{项目效果首页.png}{系统首页统计面板}
\thesisfigure{项目效果医生信息管理.png}{医生信息管理页面}
\thesisfigure{项目效果医药信息管理.png}{药品信息管理页面}
\thesisfigure{项目效果销售地点地图管理.png}{销售地点地图管理页面}
\thesisfigure{项目效果医保政策管理.png}{医保政策管理页面}
\thesisfigure{项目效果必备材料管理.png}{必备材料管理页面}
\thesisfigure{项目效果增删改查.png}{通用增删改查操作效果}

% =============================================================
\chapter{系统测试}

\section{测试策略}
测试采用分层策略：后端使用 JUnit、Mockito、Spring Boot Test 与 JaCoCo；前端使用 Vitest、Vue Test Utils 与 V8 Coverage；Python 测试验证 API 回归脚本本身；历史部署环境使用 57 项黑盒 API 用例验证容器、数据库、Redis 与 Nginx 链路。测试环境以 JDK 17、Maven 3.8+、Node 18+、MySQL 8、Redis 7 和 Docker Compose v2 为基线。

\section{代码级自动化测试结果}
\begin{table}[H]
  \centering
  \caption{代码级自动化测试结果}
  \begin{tabularx}{\textwidth}{p{3.1cm}p{2.5cm}p{2.5cm}L}
    \toprule
    测试层次 & 通过数 & 总数 & 结论 \\
    \midrule
    后端 JUnit 测试 & 100 & 100 & 全部通过 \\
    前端 Vitest 测试 & 133 & 133 & 全部通过 \\
    Python 测试 & 8 & 8 & 全部通过 \\
    合计 & 241 & 241 & 通过率 100\% \\
    \bottomrule
  \end{tabularx}
\end{table}

\section{覆盖率结果}
\begin{table}[H]
  \centering
  \caption{自动化测试覆盖率}
  \begin{tabularx}{\textwidth}{p{3.0cm}p{3.0cm}p{2.6cm}L}
    \toprule
    对象 & 指标 & 覆盖率 & 说明 \\
    \midrule
    后端 & 行覆盖率 & 100\%（935/935） & 达到行覆盖率目标 \\
    后端 & 方法覆盖率 & 100\% & 全部方法被执行 \\
    后端 & 指令覆盖率 & 99.74\% & 少量编译级指令未覆盖 \\
    后端 & 分支覆盖率 & 87.39\% & 复杂异常组合仍可补充 \\
    前端 & 语句覆盖率 & 100\% & 全部语句被执行 \\
    前端 & 行覆盖率 & 100\% & 全部代码行被执行 \\
    前端 & 分支覆盖率 & 95.64\% & 少量防御分支未命中 \\
    前端 & 函数覆盖率 & 63.88\% & 视图回调和框架绑定可提升 \\
    \bottomrule
  \end{tabularx}
\end{table}

\section{历史部署环境 API 回归}
2026 年 7 月 13 日的容器化环境完成 57 项 API 回归并全部通过，平均响应时间约 309.9 ms，P95 为 550 ms，最大值为 639 ms。该结果用于证明当时版本的部署链路完整性，与 7 月 15 日整理的 241 项代码级测试属于不同测试批次。

\begin{table}[H]
  \centering
  \caption{历史 API 回归结果}
  \begin{tabularx}{\textwidth}{p{3.1cm}p{2.6cm}p{2.6cm}L}
    \toprule
    指标 & 数值 & 状态 & 说明 \\
    \midrule
    API 用例 & 57/57 & 通过 & 覆盖认证、权限、业务查询和写操作 \\
    平均响应时间 & 309.9 ms & 正常 & 历史容器测试环境统计 \\
    P95 响应时间 & 550 ms & 正常 & 95\% 请求不超过该值 \\
    最大响应时间 & 639 ms & 正常 & 未发现超时或失败请求 \\
    MySQL/Redis & UP/UP & 通过 & 健康检查确认依赖可用 \\
    \bottomrule
  \end{tabularx}
\end{table}

\section{代码检查问题整改}
2026 年 7 月 15 日依据华为云 CodeArts 代码扫描报告对致命与严重问题集中整改，每项修复均附回归测试，并以小粒度提交同步至双仓库。整改后后端 JUnit 测试由 84 项增至 100 项，前端 Vitest 测试由 123 项增至 133 项。整改清单见表\ref{tab:fix}。

\begin{table}[H]
  \centering
  \caption{代码检查问题整改清单}\label{tab:fix}
  \begin{tabularx}{\textwidth}{p{1.1cm}p{3.6cm}L}
    \toprule
    级别 & 问题 & 整改措施 \\
    \midrule
    致命 & 临时密码用 \texttt{Math.random()}（G.OTH.01） & 改用 \texttt{SecureRandom}，补密码格式回归测试 \\
    严重 & 嵌套三元表达式（G.EXP.03） & 拆分为单层三元与中间变量 \\
    严重 & import 未分组（G.FMT.03） & 6 个生产文件补空行分组与字母排序 \\
    缺陷 & 核心业务表硬删除 & 改 \texttt{deleted\_at} 软删除，查询过滤，可恢复 \\
    缺陷 & 数据改动无追溯 & 新增 \texttt{create\_by}/\texttt{update\_by}/\texttt{deleted\_by} 审计列 \\
    缺陷 & 异常统一返回 200 & \texttt{GlobalExceptionHandler} 映射 400/401/403/404/409/500 \\
    缺陷 & 搜索每次按键发请求 & 新增 debounce，6 个管理页防抖 300\,ms \\
    加固 & 无 CSP 与定时备份 & nginx 下发 CSP；新增 mysqldump 备份脚本与 systemd timer \\
    \bottomrule
  \end{tabularx}
\end{table}

\thesisfigure{华为云测试用例.png}{CodeArts 测试用例管理}

% =============================================================
\chapter{构建部署与持续集成}

\section{本地构建与容器运行}
后端在 \texttt{medical-backend} 目录执行 \texttt{mvn verify}，生成包含测试门禁的可执行 JAR；前端在 \texttt{medical-managerment-system} 目录执行 \texttt{npm run test:coverage} 和 \texttt{npm run build}，生成 \texttt{dist} 静态资源。根目录执行 \texttt{docker compose up -d --build} 可构建并启动 \texttt{medicine-backend} 与 \texttt{medicine-web} 两个服务。

\begin{table}[H]
  \centering
  \caption{构建与运行入口}
  \begin{tabularx}{\textwidth}{p{2.8cm}p{5.2cm}L}
    \toprule
    对象 & 命令或入口 & 结果 \\
    \midrule
    后端 & \texttt{mvn verify} & 执行测试、覆盖率门禁并生成 JAR \\
    前端测试 & \texttt{npm run test:coverage} & 执行 133 项测试并生成 V8 覆盖率 \\
    前端构建 & \texttt{npm run build} & 生成生产静态资源 \texttt{dist/} \\
    Compose & \texttt{docker compose up -d --build} & 构建并启动前后端容器 \\
    健康检查 & \texttt{/actuator/health} & 返回应用、MySQL 和 Redis 状态 \\
    \bottomrule
  \end{tabularx}
\end{table}

\section{端口转发与反向代理}
开发或演示环境对外暴露 Web/Nginx 的 9092 端口，Nginx 将 \texttt{/api/}、\texttt{/image/} 和 \texttt{/actuator/health} 分别转发至后端 API、图片目录和健康检查。Compose 将后端 8082 绑定至宿主机回环地址；生产环境只开放 80/443，由外层 Nginx 统一接入。

\begin{table}[H]
  \centering
  \caption{端口与路径转发}
  \begin{tabularx}{\textwidth}{p{4.6cm}p{4.6cm}L}
    \toprule
    入口 & 目标 & 用途 \\
    \midrule
    9092 或 80/443 & \texttt{medicine-web:80} & 页面统一入口 \\
    \texttt{/api/} & \texttt{backend:\allowbreak 8082/\allowbreak api/} & 业务 API \\
    \texttt{/image/} & 后端受控上传目录 & 药品图片访问 \\
    \texttt{/\allowbreak actuator/\allowbreak health} & \texttt{backend:8082} & 健康检查 \\
    \texttt{127.0.0.1:8082} & Spring Boot & 仅限本机诊断 \\
    \bottomrule
  \end{tabularx}
\end{table}

\section{子域名与 HTTPS}
本项目正式子域名已确认为 \texttt{sx.weavecan.com}，通过 DNS 指向 ECS 并经 SSH 隧道转发至本地 Docker 部署，TLS 证书已配置。Nginx \texttt{server\_name}、后端 CORS 白名单与 Cookie 安全策略均已按该域名同步。正式访问地址为 \url{https://sx.weavecan.com/}。

\begin{table}[H]
  \centering
  \caption{生产域名上线核对项}
  \begin{tabularx}{\textwidth}{p{3.5cm}p{4.2cm}L}
    \toprule
    配置项 & 建议值 & 说明 \\
    \midrule
    DNS & \texttt{sx.weavecan.com} 指向 ECS & A、AAAA 或 CNAME \\
    Nginx \texttt{server\_name} & \texttt{sx.weavecan.com} & 替换演示主机名 \\
    CORS & \texttt{https://\allowbreak sx.weavecan.com} & Cookie 模式不能与通配符同时使用 \\
    Cookie Secure & \texttt{true} & HTTPS 生产环境必须开启 \\
    防火墙/安全组 & 22（受限）、80、443 & 关闭 8082、MySQL、Redis 公网暴露 \\
    \bottomrule
  \end{tabularx}
\end{table}

\section{CodeArts 流水线}
CodeArts Repo 使用 \texttt{shixun-2/master}。已归档的流水线 \#12 对提交 \texttt{41ade53e} 执行构建、代码检查、制品和部署，7 月 14 日全流程通过，总耗时 5 分 1 秒；部署目标为 ECS \texttt{1.14.150.130}，演示端口 9092。流水线最终通过公网健康检查判定发布结果（阶段见表\ref{tab:pipeline}）。

\thesisfigure{华为云制品仓库.png}{CodeArts 制品仓库}

% =============================================================
\chapter{个人贡献与总结}

\section{个人提交统计}
本人在项目中以账号 \texttt{anmoryli}（过程文档账号 \texttt{anmory666}）合计完成 295 次提交，涵盖项目初始化、各功能分支开发与 15 次集成合并，体现了项目经理、开发与集成者多重角色。提交按类型分布见表\ref{tab:commit-type}。

\begin{table}[H]
  \centering
  \caption{提交按类型分布（合计 295 次）}\label{tab:commit-type}
  \begin{tabularx}{\textwidth}{p{2.4cm}p{1.6cm}L}
    \toprule
    类型 & 数量 & 说明 \\
    \midrule
    feat & 106 & 新功能：主题、安全、医生、销售、策略、认证等 \\
    fix & 60 & 缺陷修复：Vue、前端、安全、SQL、医生、Nginx、高德等 \\
    docs & 55 & 文档：README、process-docs、UML、报告 \\
    style & 19 & 代码风格：G.FMT 系列规则闭环 \\
    refactor & 13 & 重构：权限拆分、缓存事件等 \\
    chore & 13 & 杂项与构建配置 \\
    test & 9 & 测试补全 \\
    ci/security/ops & 4 & 流水线、安全加固、备份脚本 \\
    功能分支集成合并 & 15 & 集成 15 个功能分支至主分支 \\
    项目初始化 & 1 & 仓库建立与首次提交 \\
    \bottomrule
  \end{tabularx}
\end{table}

\section{项目经理贡献}
作为项目经理，本人的管理工作贯穿全程，主要包括：划分 5 人角色并建立 owner 映射；建立 Epic--Feature--Story--Task--里程碑五层工作项（5/7/15/57/5）；设定 5 个里程碑的验收标准与负责人并负责 MS-01、MS-02；撰写双仓库提交规范并推动落地；配合配置 CodeArts 流水线并定义 PR 协作流程；通过 4 轮代码扫描驱动 finding 闭环并将覆盖率门禁固化为构建硬门禁；建立 \texttt{evidence/} 三级证据归档机制；承担 10 个认证黑盒验收任务。

\section{开发贡献}
作为开发工程师，本人主导 EP-01“统一身份与权限中心”，在后端 42 个 API 中承担认证、权限、上传、仪表盘与高德 \texttt{/api/regeo} 代理；主导 httpOnly Cookie 持久登录改造（后端 8 个小粒度提交）；推动 V5 规范化 RBAC 迁移，将“账号字符串角色”升级为账号--角色--权限模型；主导代码扫描致命与严重问题整改（\texttt{SecureRandom}、软删除、审计列、HTTP 状态码映射、搜索防抖等）；并完成最终交付文档与个人实习报告的撰写。

\section{代码规模}
\begin{table}[H]
  \centering
  \caption{项目代码规模}
  \begin{tabularx}{\textwidth}{p{4.0cm}p{3.0cm}L}
    \toprule
    对象 & 行数 & 说明 \\
    \midrule
    后端 Java & 6\,387 & Spring Boot 业务、安全、数据访问 \\
    前端 Vue/JS & 9\,294 & Vue 3 页面、路由、状态与 API \\
    合计 & 约 15\,681 & 不含 SQL 脚本、配置与文档 \\
    \bottomrule
  \end{tabularx}
\end{table}

\section{质量保证与风险控制}
\begin{table}[H]
  \centering
  \caption{主要风险与控制措施}
  \begin{tabularx}{\textwidth}{p{3.2cm}p{4.8cm}L}
    \toprule
    风险 & 影响 & 控制措施 \\
    \midrule
    生产密钥泄露 & 数据库/Redis 被未授权访问 & Secret 与环境变量，不提交真实密码，定期轮换 \\
    权限绕过 & 医生获得管理员写权限 & 服务端角色校验、越权测试、权限树只作界面辅助 \\
    图片路径穿越 & 任意文件被覆盖或读取 & 路径规范化、随机文件名、目录边界校验 \\
    数据迁移失败 & 部署中断或数据不一致 & 迁移前备份、版本化脚本、预发布验证与回滚 \\
    双库历史分叉 & GitHub 与 CodeArts 不一致 & 同一提交同步、合并前拉取、禁止强推 \\
    覆盖率误读 & 以行覆盖代替真实质量 & 同时报告分支、函数、API 与端到端结果 \\
    \bottomrule
  \end{tabularx}
\end{table}

\section{实习心得体会}
本次实训让我在四个工作日内完整经历了一个软件工程项目的全生命周期，收获主要体现在三方面。其一，\textbf{项目管理层面}，我体会到“先分解再执行”的价值——通过 Epic/Feature/Story/Task 分层与里程碑验收标准，团队五人能在短时间内并行推进而不混乱，双仓库规范与 CI/CD 流水线让协作有秩序、交付可追溯。其二，\textbf{工程能力层面}，规范化 RBAC、httpOnly Cookie 会话、缓存失效事件、路径安全上传等设计让我深入理解了“安全边界必须在服务端”的原则，而 100\% 行覆盖门禁与代码扫描闭环让我养成了“每改一处必跑测试、每发现一条规则必小粒度修复”的习惯。其三，\textbf{工程习惯层面}，最小粒度提交、推送前校验与“每改一处必跑测试、每发现一条规则必小粒度修复”的纪律，让我体会到工程规范对个人开发效率与代码质量的放大作用。

\section{不足与展望}
本次实习亦存在不足：本轮执行机无 Docker/MySQL/Redis/管理员凭据，未执行实时部署级 57 用例 API 回归，历史 57/57 仅作基线证据；行覆盖率 100\% 不等于所有业务组合无缺陷。后续改进包括：在隔离集成环境重新执行 57 项 API 回归并设为发布门禁；增加 Playwright/Cypress 浏览器端到端测试；提升前端函数覆盖率与前后端分支覆盖率；将 MySQL 与 Redis 迁入私网并启用 TLS；完善正式子域名的 WAF、日志采集、指标告警与备份恢复演练。

% =============================================================
\addcontentsline{toc}{chapter}{参考文献}
\begin{thebibliography}{99}
  \bibitem{springboot} VMware. Spring Boot Reference Documentation[EB/OL].
  \bibitem{springsecurity} VMware. Spring Security Reference[EB/OL].
  \bibitem{vue} Vue.js Team. Vue 3 Documentation[EB/OL].
  \bibitem{vite} Vite Team. Vite Documentation[EB/OL].
  \bibitem{docker} Docker Inc. Docker Compose Documentation[EB/OL].
  \bibitem{mysql} Oracle. MySQL 8.0 Reference Manual[EB/OL].
  \bibitem{redis} Redis Ltd. Redis Documentation[EB/OL].
  \bibitem{codearts} 华为云. CodeArts 软件开发生产线文档[EB/OL].
  \bibitem{projectrepo} 慧医数字医疗应用系统项目仓库、SQL、部署脚本和过程文档[Z]. 2026.
\end{thebibliography}

\appendix
\chapter{主要 API 清单}

\begin{longtable}{p{2.1cm}p{4.4cm}p{2.8cm}p{4.2cm}}
  \caption{认证与公共 API}\label{tab:public-api}\\
  \toprule
  方法 & 路径 & 权限 & 说明 \\
  \midrule
  \endfirsthead
  \toprule
  方法 & 路径 & 权限 & 说明 \\
  \midrule
  \endhead
  POST & \texttt{/api/login} & 公开 & 表单或 JSON 登录并设置 httpOnly Cookie \\
  GET & \texttt{/api/session} & 公开 & 探测并恢复当前会话 \\
  GET & \texttt{/api/permissions} & 已登录 & 按 accountId 返回 MENU 树、ACTION 权限码与角色并集 \\
  POST & \texttt{/api/logout} & 已登录 & 删除 Redis 会话并清除 Cookie \\
  GET & \texttt{/api/dashboard} & \texttt{dashboard:\allowbreak read} & 返回首页统计、分布和资讯 \\
  GET & \texttt{/actuator/health} & 网关策略控制 & 返回应用、MySQL 和 Redis 状态 \\
  GET & \texttt{/api/regeo} & \texttt{sale-map:\allowbreak read} & 高德逆地理编码代理 \\
  \bottomrule
\end{longtable}

\begin{longtable}{p{2.5cm}p{5.2cm}p{5.8cm}}
  \caption{业务 API 概览}\label{tab:business-api}\\
  \toprule
  模块 & 查询接口（ACTION） & 写操作接口（ACTION） \\
  \midrule
  \endfirsthead
  \toprule
  模块 & 查询接口（ACTION） & 写操作接口（ACTION） \\
  \midrule
  \endhead
  城市 & \texttt{GET /api/citys}（\texttt{city:read}） & \texttt{POST}、\texttt{DELETE /api/citys/\{id\}}（\texttt{city:write}） \\
  公司 & \texttt{GET /api/companys}（\texttt{company:read}） & \texttt{POST}、\texttt{PUT/DELETE}（\texttt{company:write}） \\
  销售地点 & \texttt{GET /api/sales}（\texttt{sale:read}） & \texttt{POST}、\texttt{PUT/DELETE}（\texttt{sale:write}） \\
  销售地图 & \texttt{GET /api/regeo}（\texttt{sale-map:read}） & -- \\
  公司政策 & \texttt{GET /api/company\_policys}（\texttt{company-policy:read}） & \texttt{POST}、\texttt{PUT/DELETE}（\texttt{company-policy:write}） \\
  医保政策 & \texttt{GET /api/medical\_policys}（\texttt{medical-policy:read}） & \texttt{POST}、\texttt{PUT/DELETE}（\texttt{medical-policy:write}） \\
  医生 & \texttt{GET /api/doctors}（\texttt{doctor:read}） & 新增/修改/删除（\texttt{doctor:write}）；密码重置（\texttt{doctor:reset-password}） \\
  药品 & \texttt{GET /api/drugs}（\texttt{drug:read}） & \texttt{POST}、\texttt{PUT/DELETE}（\texttt{drug:write}） \\
  材料 & \texttt{GET /api/materials}（\texttt{material:read}） & \texttt{POST}、\texttt{PUT/DELETE}（\texttt{material:write}） \\
  图片 & \texttt{GET /image/\{name\}}（静态资源策略） & \texttt{POST /api/base/upload}（\texttt{file:upload}） \\
  \bottomrule
\end{longtable}

\end{document}
